The ending of v5 for the \emph{v4 row} remains correct: that row
was \UNVALIDATED.
It was not ``cone anomaly found,'' and it was not ``BRST restored.''
v6 replaces it with a derived row, analysed below.

\subsection{Diagnosis of the v4 failure (new in v6)}
\label{sec:STdiag}

The frozen XVI-P antifield vertex has the closed form
\(K'_{(mn)}\,a\,(\partial_mc_n+\partial_nc_m)\), summed over the
independent coordinates \((mn)\), \(m\le n\) \cite{CheungP}.
Writing the BRST source term in the same independent coordinates
fixes three conventions that v4 left free:
\begin{enumerate}[label=\textup{(C\arabic*)}]
\item the ghost enters the vertex with an upper index, and
\(c_n=\eta_{nn}c^n\) appears in
\(\delta_ch_{(mn)}=p_mc_n+p_nc_m\);
\item the equation-of-motion factor is the derivative with respect
to the independent coordinate,
\[
\frac{\partial S}{\partial h_{(mn)}}=(2-\delta_{mn})E^{mn}(h),
\]
with \(E\) the linearized Einstein operator and indices raised by
\(\eta=\diag(1,-1,-1,-1)\);
\item consequently the \(K'_{(03)}\) coordinate carries weight
\(2E^{03}=-2E_{03}\).
In lower-index components this is weight \(w=-2\).
\end{enumerate}
\textsc{[exec]} The v4 convention scan (\loc{code/st_scan.py},
\loc{logs/v4/st_scan.json}) used only
positive weights \(w\in\{\tfrac12,1,2\}\).
It never contained \(w=-2\) on \((03)\) together with (C1).
This is why none of its 72 conventions closed.
The UNIT ``nonzero source term'' in v4 is produced by the same
weight error.

\subsection{An \emph{a priori} normalization}

v5 rightly warned that a fitted \(\lambda\) proves nothing.
In v6, \(\lambda\) is fixed before the loop data are used.
The XVI-P tree data give the antighost--ghost normalization
\(A_{\rm GH}=\tfrac12\).
They also give the quadratic-action normalization \(N=1\) of
\(E^{\rm up}\) on de Donder configurations.
\textsc{[exec]} \loc{code/st_norm.py} solves the de Donder
constraint for \((h_{00},h_{03})\), with \(h_{11}=h_{22}\) tied, and
returns
\[
\lambda_{\rm pred}=\frac{N}{\ii A_{\rm GH}}=-2\ii
\]
(\loc{logs/v6/st_norm.json}).

\subsection{Constant-profile closure on CR and F}

Let the axisymmetric lattice coordinates be
\(x=(h_{00},h_{03},h_{33},f_0)\), with
\(h_{11}=h_{22}=f_0\), and the ghost \(c=(c^0,c^3)\).
Put \(p_\mu=(q_0,0,0,q_3)\).
Define
\begin{equation}
\mathcal S_\lambda(c,h)=
\sum_{a,b}\widehat B^{hh}_{ab}\,(\delta_ch)_a\,h_b
+\lambda\sum_{(mn),\rho}
\widehat\Gamma^{\mathrm{D13}}_{(mn),\rho}\,c^\rho\,
(2-\delta_{mn})E^{mn}(h),
\label{eq:STv6}
\end{equation}
where
\(\widehat B^{hh}=\mathrm{D01+D02+D03+D11+D12}\)
denotes the executed constant-profile poles restricted to lattice
keys.

\begin{theorem}[Constant-profile \(hh/Kc\) Slavnov--Taylor row]
\label{thm:STclosure}
\textsc{[exec]}
For both CR and F and every \(1/2<\alpha<1\), the polynomial
\(\mathcal S_\lambda\) in
\((h,f_0,c,q_0,q_3)\) vanishes identically if and only if
\(\lambda=-2\ii=\lambda_{\rm pred}\).
More precisely:
\begin{enumerate}[label=\textup{(\roman*)}]
\item Sixteen monomials receive a nonzero graph contribution.
On twelve of them, the D01 and D03 contributions cancel among the
\(hh\) graphs.
D02, D11, and D12 contribute nothing to the row.
\item The remaining four monomials are
\(f_0c^0q_0^3,\ f_0c^0q_0q_3^2,\ f_0c^3q_0^2q_3,\ f_0c^3q_3^3\).
On them, the \(hh\) gauge variation has coefficients
\(\mp\ka^2(\alpha-1)/(4\pi)=\mp\ka^2V/(2\pi)\), and the antifield
term has coefficients
\(\pm\ii\ka^2(\alpha-1)/(8\pi)\).
Hence a single \(\alpha\)-independent \(\lambda=-2\ii\) cancels
all four.
\item The matter sub-row \(\mathrm{D11}+\mathrm{D12}\) is
transverse by itself.
\item The UNIT replacement \(c_S\equiv1\) makes both terms vanish
identically, a trivial closure.
So does the flat four-dimensional benchmark on the lattice
(\loc{code/flat_pole.py}).
Neither of these can determine \(\lambda\), which is why
\(\lambda\) is fixed \emph{a priori}.
\end{enumerate}
\end{theorem}

\begin{proof}[Evidence]
\loc{code/st_v6.py} forms \eqref{eq:STv6} symbolically from the
logged poles and extracts every monomial coefficient as an exact
rational function of \(\alpha\).
It solves for \(\lambda\) on the first equation and checks
\(\partial_\alpha\lambda=0\) and every residual \(=0\)
(\loc{logs/v6/st_v6_\{CR,F,UNIT,FLAT\}.json}).
The per-graph table is in \loc{logs/v6/st_v6\_stats.json}.

Section~\ref{sec:structures} gives the structural reason.
On the lattice, the gauge variation of the \(hh\) sum is
\(2(b_6+b_7)P^2f_0\,(p\cdot c)\), so
\(\mathcal S_\lambda\propto V\) on both branches.
The branch-dependent structure \(b_8\) is gauge invariant.

The validity of the flat tensor-averaging routine is checked
separately. The flat D11 covariant ratios \([6,1,-6,-1,8]\)
reproduce the frozen \texttt{P18\_D11} record
(\loc{logs/v6/flat\_vs\_p18.json}).
\end{proof}

\begin{remark}[Scope of the closure]
\label{rem:STscope}
Theorem~\ref{thm:STclosure} answers v5's slice objection
(Example~\ref{ex:slice}) for this row by direct computation, not
by appeal to shift invariance.
It is \emph{one} row, under these restrictions:
\begin{itemize}
\item the constant profile;
\item the axisymmetric lattice \(h_{11}=h_{22}\), without the
\((12)\) entry;
\item excluding the vector-type keys on which mixed-helicity
monomials were dropped (Section~\ref{sec:catalogue});
\item without the auxiliary-field terms that vanish on this
restriction;
\item and with the spectral reduction itself conditional on
Assumption~\ref{ass:reduction}.
\end{itemize}
It is not the 26 Ward-constrained combinations, and it is not a
full Slavnov--Taylor theorem on the cone.
Within that scope the honest status changes from \UNVALIDATED\
to \emph{closed}.
No anomaly is found, and no BRST restoration is claimed.
\end{remark}
